
\section{OA-06: Tavis--Cummings bright and dark states}
\source{Week 6 Tutorial, Question 1; released Week 6 tutorial solutions.}
\subsection*{Step 1: identify the coupled subspace}
For two identical resonant atoms and one cavity mode,
\[
H_I=\hbar g(aJ_++a^\dagger J_-),\qquad
J_-=\sigma_-^{(1)}+\sigma_-^{(2)}.
\]
Define
\(\ket S=(\ket{eg}+\ket{ge})/\sqrt2\) and
\(\ket A=(\ket{eg}-\ket{ge})/\sqrt2\).
Act on each state:
\[
J_-\ket S=\sqrt2\ket{gg},\qquad
J_-\ket A=0.
\]
The symmetric state is \emph{bright} with enhanced coupling
\(\sqrt2g\); the antisymmetric state is \emph{dark}.

\subsection*{Step 2: propagate a localised excitation}
Initially \(\ket{\Psi(0)}=\ket{eg;0}
=(\ket{S;0}+\ket{A;0})/\sqrt2\).
Set \(\vartheta=\sqrt2gt\). The bright component
Rabi-oscillates to \(\ket{gg;1}\) while the dark
component does not:
\[
\ket{\Psi(t)}=\frac1{\sqrt2}\left[
\cos\vartheta\ket{S;0}+\ket{A;0}
-i\sin\vartheta\ket{gg;1}\right].
\]
Return to the individual-atom basis:
\[
\boxed{\begin{aligned}
\ket{\Psi(t)}={}&\frac{1+\cos\vartheta}{2}\ket{eg;0}
+\frac{\cos\vartheta-1}{2}\ket{ge;0}\\
&-\frac{i\sin\vartheta}{\sqrt2}\ket{gg;1}.
\end{aligned}}
\]

\subsection*{Step 3: calculate observables and conserved excitation}
The cavity has exactly one photon in the last branch:
\[
\boxed{\langle n_{\rm cav}\rangle
=\tfrac12\sin^2\vartheta}.
\]
Atom 1 is excited only in the first term, atom 2
only in the second:
\[
P(e_1)=\frac{(1+\cos\vartheta)^2}{4},\qquad
P(e_2)=\frac{(1-\cos\vartheta)^2}{4}.
\]
Thus \(\langle N_{\rm at}\rangle
=P(e_1)+P(e_2)=\tfrac12(1+\cos^2\vartheta)\)
and
\[
\boxed{\langle N_{\rm at}\rangle+
\langle n_{\rm cav}\rangle=1}.
\]
At \(\vartheta=\pi\), the excitation transfers entirely
from atom 1 to atom 2 (up to phase). But
\(P(n_{\rm cav}=1)\le1/2\): the dark half
cannot populate the cavity.
\checkmethod{Writing the initial state in the bright/dark
basis is the key shortcut. Conservation of one excitation
is a useful error check at every time.}
